\subsection{The original projectors fix the actual periods}
\label{global-gram-original-projectors}

Retain the original arithmetic setting, its given
\texttt{LevelSubgroupCharacter}, and the established seed-decay
gap range.  Write $K=K_2$, $v_K=\mu_f(K)$ and
\[
 e(k)=v_K^{-1}1_K(k)\overline{\chi_2(k)},\qquad
 p(g)=\Bigl(\prod_vd_v\Bigr)
             \overline{\prod_v\langle w_v,\sigma_v(g_v)w_v\rangle}.
\]
These are exactly the finite idempotent and infinite coefficient
in the original setting, not substitute projectors.
Let $R_e$ and $R_p$ denote the original operators
\texttt{Re} and \texttt{Rph} on
$L^2(G(F)\backslash G(\mathbb A_F))$, with its original Haar measure.
The setting defines them by their weak convolution identities.
The conclusion below identifies precisely the original target
$HH_X=\operatorname{im}(R_eR_p)$.
\cite{global-gram-projector-source}

\begin{theorem}[The actual \texttt{InHH} supplier]
\label{global-gram-inhh-theorem}
The original level is measurable and compact open, with
$0<v_K<\infty$, and satisfies the original
\texttt{LevelSubgroupNormalized} condition.  For every original
column $T$, the complete Poincare period satisfies
\begin{equation}\label{global-gram-projectors-fix}
              R_pP_uT=P_uT,
              \qquad R_eP_uT=P_uT.
\end{equation}
Together with the established convergence and representative
independence, these are the entire original \texttt{InHH}
assertion.  The original $P_u$ is a continuous linear map from
its finite-dimensional column space to $HH_X$; in particular
the original \texttt{PuLinear} assertion also holds.
\end{theorem}

\begin{proof}
\emph{The literal finite level and its measure.}
Let $K_0$ be the full stabilizer of the original finite lattice,
and let $K(g)$ be the kernel of its action on $L_f/gL_f$.
The stabilizer is compact open.  For example, integral matrices
on a finite projective lattice form a compact set; preservation
of the nondegenerate Hermitian form is closed and determines
the inverse.  The finite quotients give the specified congruence
neighborhoods.  No free global integral basis is required.
The definitions give $K(g)\subseteq K\subseteq K_0$.
Indeed every defining $K_2$ difference is in $gL_f$ for an
element of $K(g)$, and one may take all its flag coefficients
equal to zero.

All vectors occurring in the original character argument $A(k)$
are in $L_f$, and the Hermitian form is integral there.
Thus for $k\in K(g)$,
$A(k)-A(1)\in g\widehat{\mathcal O}_E$.  The prescribed
finite additive character is trivial on
$\widehat{\mathcal O}_E$, so $\chi_2(k)=\chi_2(1)=1$.
The last equality follows from the given character law and
unit norm.  As $K$ is a subgroup containing the open subgroup
$K(g)$, it is open; it is also closed and hence compact in
$K_0$.  Haar measure is therefore positive and finite on $K$.
This proves exactly \texttt{LevelSubgroupNormalized}, and the
same finite-quotient description proves measurability and local
constancy of its character.\cite{global-gram-projector-source}

\emph{The finite projector.}
The established finite phase theorem gives the literal seed law
$\Phi_u(T)(y(1,k))=\chi_2(k)\Phi_u(T)(y)$ for $k\in K$.
Its absolutely convergent Poincare sum has the same law.
Consequently, pointwise for every $y$,
\[
 \int_{G_f}e(k)P_uT(y(1,k))\,dk
  =v_K^{-1}\int_K\overline{\chi_2(k)}\chi_2(k)P_uT(y)\,dk
  =P_uT(y).
\]
This uses the actual $v_K$, without probability renormalization
of any other measure.  Since $e\in L^1(G_f)$, its strong
convolution on $L^2$ is bounded.  The original weak identity
for $R_e$ identifies it with this operator, proving the second
identity in \eqref{global-gram-projectors-fix}.

\emph{The infinite projector, including the full return sum.}
Write $\Phi=\Phi_u(T)$ and set
\[
 P_{\rm abs}(y)=\sum_{\delta\in S_u(F)\backslash G(F)}
                                    |\Phi(\delta y)|.
\]
The finite column decomposition bounds
$|\Phi|\le M_TF_u^+$ for a fixed finite $M_T$, using precisely
the auxiliary positive seed of
Theorem~\ref{global-gram-framed-sum}.  Its proved positive
Poincare series lies in $L^2$; hence $P_{\rm abs}\in L^2$.
The source coefficient $p$ is in $L^1(G_\infty)$ by
Theorem~\ref{actual-seed-decay-theorem}.  Right translation
is unitary for the original quotient Haar measure.  Therefore
Minkowski's integral inequality gives
\begin{equation}\label{global-gram-convolution-majorant}
 \left\|\int_{G_\infty}|p(t)|P_{\rm abs}(\,\cdot\,(t,1))\,dt
                                                     \right\|_2
       \le\|p\|_1\|P_{\rm abs}\|_2<\infty.
\end{equation}
Tonelli, applied to this nonnegative majorant, shows that for
almost every $y$ the entire series may be interchanged with
the noncompact $t$ integral.  The original pointwise seed
reproduction then gives
\begin{align*}
 \int_{G_\infty}p(t)P_uT(y(t,1))\,dt
 &=\sum_\delta\int_{G_\infty}p(t)\Phi(\delta y(t,1))\,dt\\
 &=\sum_\delta\Phi(\delta y)=P_uT(y).
\end{align*}
In particular no compact truncation or omitted rational return
is used.  Since $p\in L^1$, the same weak-to-strong convolution
identification using the literal \texttt{IsOpInt} contract
gives $R_pP_uT=P_uT$ in $L^2$.

The earlier folding theorem supplies convergence at every point,
independence of the original rational transversal and $L^2$
membership.  The phase theorem supplies the seed invariance
and measurability.  Thus all conjuncts of the original
\texttt{InHH} statement now hold.  Linearity follows termwise
from absolute convergence; it is continuous because the source
column space is finite dimensional.  This proves the original
\texttt{PuLinear}, with no change to its map or norm.
\end{proof}

\paragraph{The exact target cut that has been removed.}
The actual $e$ and $p$ are self-adjoint convolution idempotents:
for $e$ this follows directly from the unitary character and
the displayed level volume, and for $p$ from Schur orthogonality
with its prescribed formal degrees.  Their convolution operators
commute because the finite and infinite factors commute.
Hence $Q_H=R_eR_p$ is the orthogonal projector onto $HH_X$.
For the actual normalized maps $A_u=P_u/\sqrt{c_u}$, the theorem
gives, for every pair $u,u'$ and every retained source inclusion,
\[
                  A_u^*Q_HA_{u'}=A_u^*A_{u'}.
\]
Thus the period Gram identity is an identity in precisely this
original target, with zero loss from $Q_H$.  If a further
target projector $Q$ is prescribed inside $HH_X$, its loss
$A_u^*(Q_H-Q)A_{u'}$ remains.  The original arithmetic chain
dictionary and its quantitative trace estimates are separate
obligations.
